\addchap{Preface}

This book began with a list of textbooks that should exist but do not. One entry read, in effect: textbooks on optimization end with convex analysis and with the theorem that stochastic gradient descent finds a stationary point, while what decides whether a large network trains well is something else entirely: Adam, warmup, weight decay, the learning-rate schedule, Muon, and the edge of stability. The theory of these practices exists, but it is scattered across several hundred papers, blog posts and theses, written in a dozen notations, and it has never been assembled into a book that a graduate student can learn from.

The book that resulted follows two lines. The first is a unification: almost every optimizer used to train neural networks is steepest descent under some norm, combined with moving averages and preconditioning. SGD, sign descent, Adam, Lion, Shampoo and Muon are all instances, and the norm of a whole network, assembled from the norms of its layers, tells how the learning rate must scale with width. This line was developed by Jeremy Bernstein, Laker Newhouse and their coauthors, and Newhouse's thesis at the Massachusetts Institute of Technology gives a connected account of it; a thesis, however, is not a textbook, and it does not treat the second line. The second line is about dynamics: what happens when these optimizers run with the step sizes that practice uses. Sharpness rises until the step size is barely stable and then stabilizes itself; warmup is needed because the first steps meet too much curvature; gradient noise sets a floor that a schedule removes, which explains the shapes of cosine and warmup--stable--decay schedules; weight decay drives normalized layers to an equilibrium; and the reasons that Adam beats gradient descent on transformers turn out to involve both lines at once.

Each of the fifteen chapters starts from a practice, a rule or default as it is used ``in the wild'', builds the smallest idealization in which the practice can be analysed, proves the main theorem of that idealization completely, compares its predictions with a simulation, and then says what the idealization misses and what the theory does and does not tell a practitioner to do. All chapters share one notation, summarized on the following pages, and each ends by translating the notation of its main sources into that of the book. An appendix collects the mathematics the chapters share.

The intended readers are graduate students and researchers in machine learning, optimization, statistics and applied mathematics, and practitioners who train large models and want to know which of their rules are understood, in what sense, and how far the understanding reaches. The prerequisites are undergraduate analysis, linear algebra and probability, and some acquaintance with neural networks. The prologue explains the plan of the book and works its first example.

The book is the third volume of \emph{The Missing Textbooks}. Its companion, \emph{The Hydrogen Atoms of Deep Learning}, collects solvable models of trained networks; one of them, the two-parameter product model of the edge of stability, is solved there and cited here rather than repeated.

A theory of practice has to be honest about how much of practice it explains. Some practices in this book are explained well, some partly, and some not at all, and the chapters say which. If the book succeeds, a reader will know not only why the recipes work where they do, but also exactly where the explanations stop.

\paragraph{A note on authorship.} This book was written by Claude, an AI model made by Anthropic, and conceived and edited by Xiaoyu Li, who proposed the subject, its two lines, and its position relative to existing work. The chapters were drafted from a common plan, their proofs re-derived and their numerical claims checked against simulations, and their references checked against sources wherever these could be reached; the colophon at the end of the book describes how it was made. Mistakes surely remain. Readers who find one are invited to treat it as the book treats its idealizations: as a precise statement of what is missed, and a pointer to what comes next.
